\chapter{分布式内存多重网格}
\label{ch:mpi}

\begin{keybox}
MPI 的根本变化不是“线程变成进程”，而是\textbf{一个全局状态被拆成多个 rank-owned 状态，而跨所有权读取必须显式建立数据副本与版本一致性。}因此分布式多重网格的核心链条是
\begin{equation}
\boxed{
\text{global state}
\rightarrow
\text{owned pieces}
\rightarrow
\text{ghost replicas}
\rightarrow
\text{message-driven readiness}
\rightarrow
\text{local stencil}.
}
\label{eq:mpi-mental-model}
\end{equation}
Halo exchange、global reduction 和 coarse-grid rank consolidation 都可以从这条链自然推出。
\end{keybox}

\section{全局状态与本地状态}

设全局未知量概念上为 $u$。分区以后
\begin{equation}
\Omega
=
\bigcup_{p=0}^{P-1}
\Omega_p^{owned},
\qquad
\Omega_p^{owned}
\cap
\Omega_q^{owned}
=
\varnothing.
\end{equation}

rank $p$ 真正保存的是
\begin{equation}
u_p^{local}
=
\left\{
u_p^{owned},
u_p^{ghost}
\right\}.
\label{eq:mpi-local-state}
\end{equation}

全局向量 $u$ 是这些 owned pieces 的逻辑并集，而不是某个 rank 上实际存在的一整块数组。Ghost 则不是“额外未知量”，而是\textbf{其他 rank 所有状态的本地副本}。

因此一个更准确的数据模型是：
\begin{equation}
\boxed{
\text{owned value}
=
\text{authoritative version},
\qquad
\text{ghost value}
=
\text{cached replica}.
}
\label{eq:mpi-owned-ghost-version}
\end{equation}

一旦 owner 更新了自己的 $u^{owned}$，其他 rank 上对应 ghost 就可能变旧。Halo exchange 的本质，就是把这些 replica 刷新到后继 stencil 所需的版本。

对 stencil operator，本地算子可以写成
\begin{equation}
A_p:
u_p^{local}
\longrightarrow
r_p^{owned}.
\label{eq:mpi-local-operator}
\end{equation}
只有当所需 ghost versions 已经与 owner 当前版本一致时，本地算子才与全局 stencil 对应。

\section{Ghost 数据与版本一致性}

考虑 rank $p$ 边界附近的 residual task。它需要一个由 rank $q$ 拥有的邻居值 $u_j$。从\chapref{ch:mg-parallel-operations}的角度看，这仍然只是一个普通 stencil gather：
\begin{equation}
r_i
=
b_i-(Au)_i.
\end{equation}

真正新增的是一条跨地址空间数据边：
\begin{equation}
u_j^{owned,q}
\longrightarrow
u_j^{ghost,p}
\longrightarrow
r_i^{owned,p}.
\label{eq:mpi-owner-ghost-dag}
\end{equation}

因此 MPI 通信并不是算法之外附加的一段“网络代码”，而是实现式~\eqref{eq:mpi-owner-ghost-dag}这条 producer--consumer 依赖的物理机制。

这也说明：
\begin{equation}
\boxed{
\text{message completion}
\neq
\text{整个 level 必须全局同步}.
}
\end{equation}
只有真正依赖某个 ghost 的 boundary tasks 必须等待该消息；完全本地的 interior tasks 可以更早执行。

\section{Halo 交换}
\label{sec:mpi-halo}

一次典型 halo exchange 可以分解为：
\begin{enumerate}
\item 识别哪些 owned values 是远端 rank 的输入；
\item pack 或建立可直接发送的数据视图；
\item 发布 receive/send；
\item 让通信向前推进；
\item 收到数据后写入 ghost region；
\item 标记依赖这些 ghost 的 boundary tasks 为 ready。
\end{enumerate}

如果边界数据在内存中连续，send buffer 可以直接引用原数组的一段；若边界是 strided face、复杂 patch 或非连续索引，则可能需要 pack/unpack。于是通信时间不应只写成网络本身：
\begin{equation}
T_{\mathrm{halo}}
\approx
T_{\mathrm{pack}}
+
T_{\mathrm{network}}
+
T_{\mathrm{unpack}}
+
T_{\mathrm{wait/unhidden}}.
\label{eq:mpi-halo-cost}
\end{equation}

在结构网格里，二维主要交换边带，三维主要交换 face halos。若 stencil 需要 edge/corner neighbors，可以直接交换更宽区域，也可以通过多阶段 face exchange 传播。选择哪种方式会改变 message count、message size 与 pack pattern。

\section{Interior/Boundary 分解与通信重叠}

把本地 residual tasks 分成
\begin{equation}
\mathcal T_p
=
\mathcal T_p^{interior}
\cup
\mathcal T_p^{boundary},
\end{equation}
其中 interior tasks 不依赖 remote ghost，boundary tasks 依赖通信完成。

于是自然得到
\begin{equation}
\boxed{
\text{post halo}
\rightarrow
\begin{cases}
\text{network progress},\\
\text{compute interior}
\end{cases}
\rightarrow
\text{ghost ready}
\rightarrow
\text{compute boundary}.
}
\label{eq:mpi-overlap-pattern}
\end{equation}

这比“nonblocking MPI 可以 overlap”更准确。真正产生 overlap 至少需要三个条件：

\begin{enumerate}
\item 有足够的 interior work 与通信独立；
\item MPI/runtime/NIC 确实能在这段时间推进通信；
\item pack、同步或依赖关系没有把本来可并发的两条路径重新串起来。
\end{enumerate}

如果发布 nonblocking send/recv 后立刻 wait，或者 interior 太小，
\begin{equation}
T_{\mathrm{interior}}
\ll
T_{\mathrm{network}},
\end{equation}
那么异步 API 并不会自动减少 wall time。

可以用一个理想化模型表示可隐藏部分：
\begin{equation}
T_{\mathrm{overlap}}
\approx
\max
\left(
T_{\mathrm{interior}},
T_{\mathrm{network}}
\right)
+
T_{\mathrm{boundary}}
+
T_{\mathrm{unhidden\ overhead}}.
\end{equation}

\section{Surface-to-Volume 与 Strong Scaling}

对边长约为 $n$ 的三维立方子域，一层六面 halo 约含
\begin{equation}
6n^2
\end{equation}
个 values，而本地 unknown 数约为
\begin{equation}
n^3.
\end{equation}
因此
\begin{equation}
\boxed{
\frac{\text{halo values}}
{\text{owned cells}}
\sim
\frac{6}{n}.
}
\label{eq:mpi-surface-volume-ratio}
\end{equation}

固定全局问题做 strong scaling 时，rank 数增加会让每 rank 的局部 $n$ 下降，于是通信/计算比自然上升。这一结论与网络多快无关，是几何分解本身造成的。

若每方向把局部长度减半，owned work 约减少到原来的 $1/8$，而 face halo 只减少到原来的 $1/4$。因此通信相对占比会继续增加。

对 $d$ 维超立方子域，一般有
\begin{equation}
\frac{\text{boundary measure}}
{\text{volume}}
=
\Theta\left(\frac1n\right).
\end{equation}

这也是\chapref{ch:benchmarking}中 strong-scaling efficiency 最终下降的一个基础原因。

\section{消息数量与消息体积}

点到点消息常用
\begin{equation}
T_{\mathrm{msg}}
\approx
\alpha+\beta m
\end{equation}
描述。若一个阶段发送 $k$ 个不能完全重叠的小消息，总成本会更接近
\begin{equation}
k\alpha+\beta M,
\qquad
M=\sum_{j=1}^k m_j.
\end{equation}

因此下面两种 halo 即使总字节数相同，也可能有完全不同的时间：
\begin{itemize}
\item 少量大消息；
\item 大量很小的 patch messages。
\end{itemize}

对 block-structured AMR，这一点尤其重要。把 level 切成更多小 boxes 可以改善负载平衡，却会增加邻居关系、message count 与 metadata。于是 box granularity 同时作用于
\begin{equation}
\boxed{
\text{load balance}
\quad\leftrightarrow\quad
\text{message aggregation}.
}
\end{equation}

\section{分区拓扑与网络拓扑}

数学上，rank 邻接图由空间分区决定：谁与谁共享 face/edge/corner。物理机器上，这些 ranks 又被映射到 node、socket、NIC 与 network topology。

因此一个 halo edge
\begin{equation}
p\leftrightarrow q
\end{equation}
可能对应：
\begin{itemize}
\item 同节点共享内存 copy；
\item 跨 socket 的本地通信；
\item 同交换机下的网络通信；
\item 更远的 network route。
\end{itemize}

逻辑通信图相同，不代表物理通信成本相同。大规模运行时 rank mapping、process placement 和 topology awareness 会影响实际 latency 与 bandwidth。

本章先保留这个层次区分；具体系统配置仍以 profiling 与 benchmark 为准。

\section{MPI 多重网格操作}

\subsection{Jacobi 与 Residual}

Jacobi/residual 的 local logical task 与单进程完全相同。区别只在于 boundary tasks 的输入 readiness。

典型流程为：
\begin{enumerate}
\item 发布当前状态版本的 halo；
\item 计算 interior；
\item 等待必需 ghost ready；
\item 计算 boundary；
\item 完成当前阶段的输出版本。
\end{enumerate}

若 smoother 连续执行多次 sweep，每次会产生新的状态版本，因此通常需要新的 halo readiness；不能把上一 sweep 的 ghost 当作自动仍然有效。

\subsection{Gauss--Seidel 与着色}

RBGS 在 MPI 中同时拥有两类顺序：
\begin{equation}
\text{red}
\rightarrow
\text{black}
\end{equation}
的颜色依赖，以及跨 rank boundary 上的 ghost version 依赖。

因此一个颜色阶段可能需要：
\begin{equation}
\text{local color update}
\rightarrow
\text{halo refresh}
\rightarrow
\text{next color boundary update}.
\end{equation}

这解释了为什么同步频繁的 smoother 在 distributed memory 上可能比 Jacobi/Chebyshev 更难扩展：不是因为 MPI“不支持 GS”，而是算法 DAG 本身产生了更多跨所有权 RAW edges。

\subsection{Restriction 与 Prolongation}

如果 fine/coarse partition 保持嵌套，restriction/prolongation 可以主要在 rank 内完成。若 coarse ownership 改变，则会出现额外的 redistribution：
\begin{equation}
\text{fine owners}
\rightarrow
\text{coarse owners}.
\end{equation}

必须把两个成本分开：
\begin{equation}
\boxed{
T_{\mathrm{restriction}}
\neq
T_{\mathrm{redistribution}}.
}
\end{equation}
前者是数值 transfer kernel，后者是 ownership remapping。

\subsection{Norm 与 Krylov Dot Product}

本地先形成 partial sum
\begin{equation}
s_p
=
\sum_{i\in\Omega_p^{owned}}r_i^2,
\end{equation}
再通过 collective 得到
\begin{equation}
\|r\|_2^2
=
\sum_p s_p.
\end{equation}

邻居 halo 只依赖局部 rank graph；global reduction 则需要把所有 ranks 的贡献合并。它在依赖图中具有更广的 fan-in，因此更容易进入关键路径。

需要精确理解一点：collective 不一定在实现上等同于“所有 rank 同时停在 barrier”，但\textbf{任何后继操作若需要最终归约结果，就必须等待这条全局数据依赖完成。}

\section{Global Reduction 与扩展性}

如果一次 Krylov iteration 含多个 dot products，而每个 dot 都需要 global collective，那么即使 local stencil 已经高度优化，外层仍可能由 reduction latency 限制。

一个非常粗略的树形 collective 心智模型是
\begin{equation}
T_{\mathrm{reduce}}
\sim
\alpha \log P
+
\beta m \log P,
\end{equation}
实际算法可能使用不同拓扑、分段和硬件 offload，因此这不是通用性能公式；它只说明：\textbf{collective critical path 不会像 local work 那样随每 rank unknown 数下降。}

这正是 strong scaling 后期常见的结构：
\begin{equation}
T_{\mathrm{local\ work}}
\downarrow,
\qquad
T_{\mathrm{collective}}
\not\downarrow\ \text{at the same rate}.
\end{equation}

因此减少不必要 residual norms、选择 communication-avoiding/pipelined Krylov 或合并可合并的 reductions，都是在修改外层 DAG，而不是单纯优化 MPI 函数调用。

\section{分布式 V-Cycle 的依赖结构}

一个 level 上的典型时间线可以写成：

\begin{enumerate}
\item 为当前 smoother/operator state 发布 halo；
\item interior work 与网络尽可能并行；
\item ghost ready 后完成 boundary work；
\item residual ready；
\item restriction；
\item 如 coarse ownership 改变则 redistribution；
\item 在 coarse active rank set 上递归；
\item coarse error 返回后 prolongation/correction；
\item 必要时把 correction redistribution 回 fine owners；
\item post-smooth；
\item 只有算法需要时才做 global reduction。
\end{enumerate}

因此“每层一次 halo”不是正确的性能模型。一次 V-cycle 内可能有多次 smoother halos、operator halos、transfer redistribution 和全局 collectives。

更好的表达是：
\begin{equation}
\boxed{
T_{\mathrm{V}}
=
\text{local kernels}
+
\text{neighbor communication}
+
\text{global collectives}
+
\text{ownership redistribution}
+
\text{unhidden waits}.
}
\end{equation}

这些项之间可能重叠，最终 wall time 由 DAG critical path 决定，而不是简单字节数相加。

\section{粗化与 Rank 并行度衰减}

由\chapref{ch:mg-parallel-operations}，
\begin{equation}
N_{\ell+1}
\approx
\frac{N_\ell}{2^d}.
\end{equation}
如果一直保留全部 $P$ 个 ranks，
\begin{equation}
\frac{N_\ell}{P}
\end{equation}
会快速下降。

当每 rank 只剩很少 unknowns 时，会同时出现：
\begin{itemize}
\item local kernel 太短；
\item halo messages 变小，latency 比例上升；
\item 每 rank box 数太少，load balance 粒度变差；
\item collective 仍然涉及过多 ranks；
\item bottom solver 需要处理大量几乎空闲的 participants。
\end{itemize}

这就是 coarse-grid rank parallelism collapse。

\section{Rank Consolidation 与所有权重映射}

设第 $\ell$ 层有 $N_\ell$ 个 unknowns，希望每个 active rank 至少拥有 $n_{\min}$ 个 unknowns，可以选择
\begin{equation}
P_\ell
\le
\min
\left(
P_{\ell-1},
\frac{N_\ell}{n_{\min}}
\right).
\label{eq:mpi-active-ranks}
\end{equation}

当
\begin{equation}
P_\ell<P_{\ell-1},
\end{equation}
coarse state 需要从旧 owner set 迁移到新的更小 rank set。这个过程可以抽象为
\begin{equation}
\text{old ownership}
\rightarrow
\text{redistribution}
\rightarrow
\text{new ownership}.
\end{equation}

它不是“关闭一些 MPI 进程”，而是\textbf{让粗层计算只在更小的 active communicator / rank subset 上执行}。

是否值得 consolidation，需要比较
\begin{equation}
\boxed{
T_{\mathrm{redistribute}}
+
T_{\mathrm{coarse}}(P_\ell)
<
T_{\mathrm{coarse}}(P_{\ell-1}).
}
\label{eq:mpi-consolidation-break-even}
\end{equation}

若 hierarchy/setup 可以复用多次 solve，一些 redistribution plan 与 metadata 可以摊销；若每次 solve 都重建，则 setup cost 也必须进入 break-even。

\section{Strong Scaling 与 Weak Scaling}

Strong scaling 固定全局问题。随着 rank 数增加：
\begin{itemize}
\item 每 rank local work 减少；
\item surface/volume ratio 上升；
\item small-message latency 相对变重；
\item coarse levels 更早出现 rank scarcity；
\item global reduction 占比通常上升。
\end{itemize}

Weak scaling 则保持每 rank local work 近似固定。此时 neighbor halo 的 surface/volume 可能近似稳定，但：
\begin{itemize}
\item global rank 数增加；
\item collective critical path 仍可能增长；
\item network contention 与 topology effects 可能增强；
\item coarse-grid communicator 仍可能需要重新设计。
\end{itemize}

因此“weak scaling 好”并不意味着 distributed solver 已经没有通信问题；它只说明局部 workload 没有像 strong scaling 那样被持续切碎。

\section{物理边界、分区边界与 AMR 接口}

一个 rank 子域边界不等于 PDE 物理边界。三类接口是：

\begin{equation}
\boxed{
\text{physical boundary}
\neq
\text{MPI partition boundary}
\neq
\text{AMR coarse--fine interface}.
}
\label{eq:mpi-boundary-classes}
\end{equation}

物理边界由\chapref{ch:discretization}中的 Dirichlet/Neumann/Robin/DtN 等关系闭合；partition boundary 只是同一 PDE 状态被不同 rank 所有；AMR coarse--fine interface 则需要层间插值、平均或通量匹配。

三者都可能使用 ghost-like storage，但数学来源完全不同。实现若只用一个“boundary flag”把它们混在一起，会同时伤害正确性与可维护性。

\section{MPI 性能诊断}

当 distributed V-cycle 扩展性变差时，建议按下面顺序定位：

\begin{enumerate}
\item \textbf{local work}：每 rank 当前还有多少 cells/tiles，kernel 是否已经太短；
\item \textbf{surface/volume}：halo bytes 相对 local work 是否快速上升；
\item \textbf{message count}：是否有大量小 patch messages，被 $\alpha$ 主导；
\item \textbf{pack/unpack}：CPU memory copy 是否已成为 halo 的主要成本；
\item \textbf{overlap}：interior work 是否真的覆盖了网络时间；
\item \textbf{progress}：nonblocking communication 是否实际向前推进；
\item \textbf{load balance}：最慢 rank 是否明显拖长阶段尾部；
\item \textbf{collective}：Allreduce 等全局依赖占比多少；
\item \textbf{coarse ownership}：是否仍让过多 ranks 活跃在极小 coarse levels；
\item \textbf{topology}：逻辑邻居与物理网络放置是否严重错配。
\end{enumerate}

如果只看“MPI 时间占比高”，仍然无法判断真正该做的是合并消息、增大 local work、改善 pack、做 overlap、减少 collectives，还是直接收缩 coarse ranks。

\section{本章小结}

分布式内存多重网格的正确心智模型不是“每个 rank 算一块，然后互相通信”，而是\textbf{全局状态被拆成 owned versions，远端读取通过 ghost replicas 实现，message completion 负责建立版本 readiness}。

Neighbor halo、global reduction 与 rank consolidation 分别对应三类不同的数据依赖：局部跨所有权读取、全局聚合、层级所有权重映射。Strong scaling 过程中 local work 下降得比 latency、collective 和固定成本更快，因此 coarse levels 最终必须改变 active rank set，而不是机械保留 finest-level 分区。

\section*{习题}

\subsection*{Owned/Ghost 与 Halo 推导}
\begin{enumerate}[label=\arabic*.]
\item 给定一维三点 stencil，把 rank $p$ 最右侧 owned cell 的 residual 写成 owner--ghost--consumer DAG，标出 authoritative value 与 replica。
\item 证明 ghost 更新完成后，本地算子式~\eqref{eq:mpi-local-operator}与对应全局 stencil 行一致。
\item 三维全局 $512^3$ 网格按 $8\times8\times8$ ranks 均分。计算每 rank owned 尺寸、一层六面 halo values 和式~\eqref{eq:mpi-surface-volume-ratio}。
\item 对同一全局问题改为 $16^3$ ranks，重复计算并比较 strong scaling 下 surface/volume 的变化。
\item 推导 $d$ 维超立方子域的 boundary/volume 数量级。
\end{enumerate}

\subsection*{通信成本与 Overlap}
\begin{enumerate}[label=\arabic*.]
\item 取 $\alpha=1.5\,\mu s$、带宽 $25\,\mathrm{GB/s}$，分别计算发送一个 $64^2$ double face 与 $16^2$ double face 的理想消息时间。
\item 比较“一个 1 MB 消息”和“1024 个总量为 1 MB 的消息”的 $\alpha$--$\beta$ 模型时间。
\item 假设 pack 为 $5\,\mu s$、network 为 $20\,\mu s$、unpack 为 $4\,\mu s$、interior work 为 $15\,\mu s$。估算理想 overlap 后的阶段时间，并指出还剩多少网络时间无法隐藏。
\item 使用 nonblocking communication 实现式~\eqref{eq:mpi-overlap-pattern}，用 timeline 证明通信是否真的与 interior work 重叠。
\item 构造一个 nonblocking MPI 几乎没有加速的案例，并区分“没有独立 work”和“没有通信 progress”两种根因。
\end{enumerate}

\subsection*{Collective 与 Krylov}
\begin{enumerate}[label=\arabic*.]
\item 解释为什么 neighbor halo 的依赖范围与 Allreduce 不同，即使二者传输总字节数相近。
\item 对一个每 iteration 含 3 次 global reductions 的 Krylov solver，记录 rank 数增加时 local SpMV、halo 与 Allreduce 的时间变化。
\item 比较“每层都计算 residual norm”的 V-cycle 与固定 cycle，统计额外 global reductions。
\item 说明 collective 并不必然是显式 barrier，但为什么需要 reduction 结果的后继 task 仍然受全局关键路径限制。
\end{enumerate}

\subsection*{Coarse Rank Consolidation}
\begin{enumerate}[label=\arabic*.]
\item 给定 $256^3\rightarrow128^3\rightarrow64^3\rightarrow32^3$ hierarchy，每 active rank 至少要求 20000 unknowns，计算式~\eqref{eq:mpi-active-ranks}给出的最大 rank 数。
\item 设 redistribution 代价为 $0.8$ ms，保留 128 ranks 的 coarse solve 为 $2.4$ ms，收缩到 16 ranks 后 coarse solve 为 $0.9$ ms。判断是否满足式~\eqref{eq:mpi-consolidation-break-even}。
\item 设计两级 active communicator，并记录 restriction redistribution 与 prolongation return path 的 owner mapping。
\item 讨论为什么 coarse rank consolidation 与 finest-level domain decomposition 不是同一个优化问题。
\end{enumerate}

\subsection*{端到端扩展性诊断}
\begin{enumerate}[label=\arabic*.]
\item 对一个 strong-scaling 失速案例，分别用 local work、surface/volume、message count、pack/unpack、load imbalance、Allreduce 与 coarse ranks 七个指标建立诊断表。
\item 为一个 V-cycle 画出 MPI trace，标记 interior compute、halo、collective 和 redistribution，并指出 critical path。
\item 设计一个数据结构，显式区分 physical ghost、MPI ghost 与 AMR coarse--fine ghost。
\item 比较纯 MPI 与较少 ranks + rank 内线程两种布局时，哪些成本会减少，哪些新成本会出现；不要只比较 rank 数。
\end{enumerate}
